% ===========================================================================================
%                                       CASOS DE USO
% ===========================================================================================

\section{Casos de uso}

Los casos de uso describen las interacciones entre los actores y el Sistema Integral de Producción, especificando las funcionalidades que el sistema proporciona para cumplir con los requerimientos funcionales definidos. Cada caso de uso representa una secuencia de acciones orientadas a alcanzar un objetivo específico dentro del sistema.

% ===========================================================================================
%                                       CU-01
% ===========================================================================================

\textbf{Caso de Inicio de sesión}

\begin{table}[H]
    \centering
    \caption{Caso de uso - Inicio de sesión}
    \label{tab:cu01}
    \begin{tabularx}{\textwidth}{|X|X|}
        \hline
        Identificador             & CU-01                                                                        \\ \hline
        Nombre                    & Inicio de sesión                                                             \\ \hline
        Objetivo                  & Permitir el acceso seguro al sistema mediante credenciales válidas.          \\ \hline
        Actor principal           & Usuario                                                                      \\ \hline
        Precondiciones            & El usuario debe estar previamente registrado en el sistema.                  \\ \hline
        Postcondiciones           & El usuario accede al sistema de acuerdo con los permisos asignados a su rol. \\ \hline
        Requerimiento relacionado & RF-01                                                                        \\ \hline
    \end{tabularx}
\end{table}

\textbf{Flujo principal}

\begin{enumerate}
    \item El usuario accede a la pantalla de inicio de sesión.
    \item El sistema solicita el nombre de usuario y contraseña.
    \item El usuario ingresa sus credenciales.
    \item El sistema valida la información proporcionada.
    \item El sistema concede acceso según el rol asignado.
\end{enumerate}

\textbf{Flujos alternos}

\begin{itemize}
    \item Credenciales inválidas:
          \begin{enumerate}
              \item El sistema muestra un mensaje de error.
              \item El usuario puede intentar nuevamente.
          \end{enumerate}
\end{itemize}

\textbf{Diagrama de caso de uso}

\begin{figure}[H]
    \centering
    \includegraphics[width=0.5\linewidth]{inicio_sesion}
    \caption{Flujo del inicio de sesión.}
\end{figure}

% ===========================================================================================
%                                       CU-02
% ===========================================================================================

\textbf{Caso de uso: Gestión de usuarios}

\begin{table}[H]
    \centering
    \caption{Caso de uso - Gestión de usuarios}
    \label{tab:cu02}
    \begin{tabularx}{\textwidth}{|X|X|}
        \hline
        Identificador             & CU-02                                                                 \\ \hline
        Nombre                    & Gestión de usuarios                                                   \\ \hline
        Objetivo                  & Administrar las cuentas de usuario y sus permisos dentro del sistema. \\ \hline
        Actor principal           & Administrador, Supervisor                                             \\ \hline
        Precondiciones            & El usuario debe haber iniciado sesión con permisos válidos.           \\ \hline
        Postcondiciones           & La información de los usuarios queda actualizada en el sistema.       \\ \hline
        Requerimiento relacionado & RF-02                                                                 \\ \hline
    \end{tabularx}
\end{table}

\textbf{Flujo principal}

\begin{enumerate}
    \item El actor accede al módulo de gestión de usuarios.
    \item El sistema muestra la lista de usuarios registrados.
    \item El actor selecciona crear, modificar o eliminar usuario.
    \item El sistema solicita la información correspondiente.
    \item El actor registra o modifica los datos requeridos.
    \item El sistema valida y almacena la información.
\end{enumerate}

\textbf{Flujos alternos}

\begin{itemize}
    \item Datos inválidos:
          \begin{enumerate}
              \item El sistema muestra un mensaje de error.
              \item El actor corrige la información ingresada.
          \end{enumerate}
\end{itemize}

\textbf{Diagrama de caso de uso}

\begin{figure}[H]
    \centering
    \includegraphics[width=0.5\linewidth]{gestion_usuarios}
    \caption{Flujo del gestión de usuarios.}
\end{figure}

% ===========================================================================================
%                                       CU-03
% ===========================================================================================

\textbf{Caso de uso: Generar reporte de producción}

\begin{table}[H]
    \centering
    \caption{Caso de uso - Generar reporte de producción}
    \label{tab:cu03}
    \begin{tabularx}{\textwidth}{|X|X|}
        \hline
        Identificador             & CU-03                                                          \\ \hline
        Nombre                    & Generar reporte de producción                                  \\ \hline
        Objetivo                  & Registrar información relacionada con la producción realizada. \\ \hline
        Actor principal           & Operador                                                       \\ \hline
        Precondiciones            & El usuario debe haber iniciado sesión en el sistema.           \\ \hline
        Postcondiciones           & El reporte de producción queda almacenado en el sistema.       \\ \hline
        Requerimiento relacionado & RF-03                                                          \\ \hline
    \end{tabularx}
\end{table}

\textbf{Flujo principal}

\begin{enumerate}
    \item El operador accede al módulo de reportes de producción.
    \item El sistema muestra el formulario de captura.
    \item El operador registra la información requerida.
    \item El sistema valida los datos ingresados.
    \item El sistema almacena el reporte de producción.
\end{enumerate}

\textbf{Diagrama de caso de uso}

\begin{figure}[H]
    \centering
    \includegraphics[width=0.5\linewidth]{reporte_prod.png}
    \caption{Flujo de reporte de producción.}
\end{figure}

% ===========================================================================================
%                                       CU-04
% ===========================================================================================

\textbf{Caso de uso: Gestión de registros de producción}

\begin{table}[H]
    \centering
    \caption{Caso de uso - Gestión de registros de producción}
    \label{tab:cu04}
    \begin{tabularx}{\textwidth}{|X|X|}
        \hline
        Identificador             & CU-04                                                            \\ \hline
        Nombre                    & Gestión de registros de producción                               \\ \hline
        Objetivo                  & Modificar o eliminar registros de producción existentes.         \\ \hline
        Actor principal           & Operador, Supervisor, Administrador                              \\ \hline
        Precondiciones            & Deben existir registros de producción almacenados en el sistema. \\ \hline
        Postcondiciones           & Los registros son actualizados o eliminados correctamente.       \\ \hline
        Requerimiento relacionado & RF-04                                                            \\ \hline
    \end{tabularx}
\end{table}

\textbf{Flujo principal}

\begin{enumerate}
    \item El actor accede al módulo de registros de producción.
    \item El sistema muestra los registros disponibles según el rol del usuario.
    \item El actor selecciona un registro.
    \item El actor decide modificar o eliminar el registro.
    \item El sistema procesa la operación solicitada.
    \item El sistema actualiza la información almacenada.
\end{enumerate}

\textbf{Reglas de negocio}

\begin{itemize}
    \item Los operadores únicamente pueden modificar o eliminar sus propios registros.
    \item Los supervisores y administradores pueden modificar o eliminar cualquier registro.
\end{itemize}

\textbf{Diagrama de caso de uso}

\begin{figure}[H]
    \centering
    \includegraphics[width=0.5\linewidth]{gestion_reg}
    \caption{Flujo de la gestión de registros de producción.}
\end{figure}


% ===========================================================================================
%                                       CU-05
% ===========================================================================================

\textbf{Caso de uso: Registro de almacén}

\begin{table}[H]
    \centering
    \caption{Caso de uso - Registro de almacén}
    \label{tab:cu05}
    \begin{tabularx}{\textwidth}{|X|X|}
        \hline
        Identificador             & CU-05                                                   \\ \hline
        Nombre                    & Registro de almacén                                     \\ \hline
        Objetivo                  & Gestionar la información de materias primas en almacén. \\ \hline
        Actor principal           & Supervisor                                              \\ \hline
        Precondiciones            & El supervisor debe haber iniciado sesión en el sistema. \\ \hline
        Postcondiciones           & La información del almacén queda actualizada.           \\ \hline
        Requerimiento relacionado & RF-05                                                   \\ \hline
    \end{tabularx}
\end{table}

\textbf{Flujo principal}

\begin{enumerate}
    \item El supervisor accede al módulo de almacén.
    \item El sistema muestra los registros existentes.
    \item El supervisor crea, modifica o elimina registros.
    \item El sistema valida la información ingresada.
    \item El sistema almacena los cambios realizados.
\end{enumerate}

\textbf{Diagrama de caso de uso}

\begin{figure}[H]
    \centering
    \includegraphics[width=0.5\linewidth]{gestion_alma}
    \caption{Flujo del registro de alamacén.}
\end{figure}


% ===========================================================================================
%                                       CU-06
% ===========================================================================================

\textbf{Caso de uso: Visualización del stock en almacén}

\begin{table}[H]
    \centering
    \caption{Caso de uso - Visualización del stock en almacén}
    \label{tab:cu06}
    \begin{tabularx}{\textwidth}{|X|X|}
        \hline
        Identificador             & CU-06                                                         \\ \hline
        Nombre                    & Visualización del stock en almacén                            \\ \hline
        Objetivo                  & Consultar el stock disponible de materias primas y productos. \\ \hline
        Actor principal           & Operador, Supervisor, Administrador                           \\ \hline
        Precondiciones            & Deben existir registros de almacén en el sistema.             \\ \hline
        Postcondiciones           & El usuario visualiza la información solicitada.               \\ \hline
        Requerimiento relacionado & RF-06                                                         \\ \hline
    \end{tabularx}
\end{table}

\textbf{Flujo principal}

\begin{enumerate}
    \item El usuario accede al módulo de almacén.
    \item El sistema muestra el stock disponible.
    \item El usuario consulta la información requerida.
\end{enumerate}

\textbf{Diagrama de caso de uso}

\begin{figure}[H]
    \centering
    \includegraphics[width=0.5\linewidth]{consulta_stock}
    \caption{Flujo de la visualización de stock.}
\end{figure}

% ===========================================================================================
%                                       CU-07
% ===========================================================================================

\textbf{Caso de uso: Visualización de la producción}

\begin{table}[H]
    \centering
    \caption{Caso de uso - Visualización de la producción}
    \label{tab:cu07}
    \begin{tabularx}{\textwidth}{|X|X|}
        \hline
        Identificador             & CU-07                                                \\ \hline
        Nombre                    & Visualización de la producción                       \\ \hline
        Objetivo                  & Consultar la producción registrada por día y turno.  \\ \hline
        Actor principal           & Supervisor                                           \\ \hline
        Precondiciones            & Deben existir registros de producción en el sistema. \\ \hline
        Postcondiciones           & El supervisor visualiza la información seleccionada. \\ \hline
        Requerimiento relacionado & RF-07                                                \\ \hline
    \end{tabularx}
\end{table}

\textbf{Flujo principal}

\begin{enumerate}
    \item El supervisor accede al módulo de producción.
    \item El sistema solicita la fecha y turno.
    \item El supervisor selecciona los parámetros de consulta.
    \item El sistema muestra la información correspondiente.
\end{enumerate}

\textbf{Diagrama de caso de uso}

\begin{figure}[H]
    \centering
    \includegraphics[width=0.5\linewidth]{ver_prod}
    \caption{Flujo de la visualización de producción.}
\end{figure}

% ===========================================================================================
%                                       CU-08
% ===========================================================================================

\textbf{Caso de uso: Programación de órdenes de producción}

\begin{table}[H]
    \centering
    \caption{Caso de uso - Programación de órdenes de producción}
    \label{tab:cu08}
    \begin{tabularx}{\textwidth}{|X|X|}
        \hline
        Identificador             & CU-08                                                        \\ \hline
        Nombre                    & Programación de órdenes de producción                        \\ \hline
        Objetivo                  & Gestionar la programación de órdenes de producción.          \\ \hline
        Actor principal           & Programador                                                  \\ \hline
        Precondiciones            & El programador debe haber iniciado sesión en el sistema.     \\ \hline
        Postcondiciones           & Las órdenes de producción quedan registradas o actualizadas. \\ \hline
        Requerimiento relacionado & RF-08                                                        \\ \hline
    \end{tabularx}
\end{table}

\textbf{Flujo principal}

\begin{enumerate}
    \item El programador accede al módulo de programación.
    \item El sistema muestra las órdenes existentes.
    \item El programador registra o modifica órdenes de producción.
    \item El sistema valida la información.
    \item El sistema almacena la programación correspondiente.
\end{enumerate}

\textbf{Diagrama de caso de uso}

\begin{figure}[H]
    \centering
    \includegraphics[width=0.5\linewidth]{prog_ordenes}
    \caption{Flujo de la programación de ordenes.}
\end{figure}

% ===========================================================================================
%                                       CU-09
% ===========================================================================================

\textbf{Caso de uso: Modificación de productos}

\begin{table}[H]
    \centering
    \caption{Caso de uso - Modificación de productos}
    \label{tab:cu09}
    \begin{tabularx}{\textwidth}{|X|X|}
        \hline
        Identificador             & CU-09                                                      \\ \hline
        Nombre                    & Modificación de productos                                  \\ \hline
        Objetivo                  & Actualizar las especificaciones de productos registrados.  \\ \hline
        Actor principal           & Administrador                                              \\ \hline
        Precondiciones            & Deben existir productos registrados en el sistema.         \\ \hline
        Postcondiciones           & Las especificaciones de los productos quedan actualizadas. \\ \hline
        Requerimiento relacionado & RF-09                                                      \\ \hline
    \end{tabularx}
\end{table}

\textbf{Flujo principal}

\begin{enumerate}
    \item El administrador accede al módulo de productos.
    \item El sistema muestra las especificaciones registradas.
    \item El administrador modifica la información requerida.
    \item El sistema valida los cambios realizados.
    \item El sistema actualiza las especificaciones correspondientes.
\end{enumerate}

\textbf{Diagrama de caso de uso}

\begin{figure}[H]
    \centering
    \includegraphics[width=0.5\linewidth]{mod_prod}
    \caption{Flujo de modificación de productos.}
\end{figure}